\documentclass[10pt,twocolumn,letterpaper]{article}

% --- Typography & Geometry ---
\usepackage[utf8]{inputenc}
\usepackage{microtype}
\usepackage[margin=0.75in,top=0.85in,bottom=0.85in]{geometry}

% --- Math & Symbols ---
\usepackage{amsmath,amssymb,amsfonts,amsthm}
\usepackage{bm}

% --- Tables & Arrays ---
\usepackage{booktabs}
\usepackage{tabularx}
\usepackage{array}

% --- Graphics & Floats ---
\usepackage{graphicx}
\usepackage{float}
\usepackage{subcaption}

% --- Code Listings ---
\usepackage{xcolor}
\usepackage{listings}

\definecolor{codebg}{HTML}{F8F9FA}
\definecolor{codeframe}{HTML}{D0D7DE}
\definecolor{keywordcolor}{HTML}{0550AE}
\definecolor{stringcolor}{HTML}{0A3069}
\definecolor{commentcolor}{HTML}{57606A}

\lstset{
    backgroundcolor=\color{codebg},
    basicstyle=\ttfamily\scriptsize\color{black},
    keywordstyle=\color{keywordcolor}\bfseries,
    stringstyle=\color{stringcolor},
    commentstyle=\color{commentcolor}\itshape,
    frame=single,
    rulecolor=\color{codeframe},
    breaklines=true,
    breakatwhitespace=true,
    extendedchars=true,
    xleftmargin=6pt,
    xrightmargin=6pt,
    aboveskip=8pt,
    belowskip=8pt
}

% --- Hyperlinks ---
\usepackage{hyperref}
\hypersetup{
    colorlinks=true,
    linkcolor=blue!70!black,
    citecolor=blue!70!black,
    urlcolor=blue!70!black
}

% --- Input Title & Authors ---
\section{00 Meta}

Draft content for vspec.


\begin{document}

\maketitle

% --- Abstract & Keywords Block ---
\begin{abstract}
Attention projection matrices in transformers are typically learned densely, 
yet their spectra are often highly non-uniform. We propose \textbf{VSpec++}, a 
variational spectral attention module that combines **three orthogonal 
gating mechanisms**:

\begin{enumerate}
\item \textbf{Token routing} -- select top-$k$ tokens via a learned saliency network
  with guaranteed \texttt{[CLS]} retention.
\item \textbf{Variational spectral gating} -- modulate each truncated singular value
  with a stochastic, input-dependent gate trained via the reparameterization
  trick.
\item \textbf{Head gating} -- weight attention heads per-example via a masked-pooled
  network.

\end{enumerate}
The Q/K/V/O projections are factorized via truncated SVD, initialized 
zero-shot from the pretrained weights, and augmented with LoRA adapters. 
All base weights (SVD factors) are frozen; only the gates, router, and LoRA 
adapters are trained.

We evaluate on DistilBERT, RoBERTa-base, and GPT-2 on SST-2 with 
{{RECOVERY\_STEPS}} recovery steps. VSpec++ achieves \textbf{{{D\_ACC}}\%} on 
DistilBERT (baseline {{BASE\_D}}\%), \textbf{{{R\_ACC}}\%} on RoBERTa (baseline 
{{BASE\_R}}\%), and \textbf{{{G\_ACC}}\%} on GPT-2 (baseline {{BASE\_G}}\%). Attention 
parameters reduce by \textbf{{{D\_RATIO}}$\times$--{{G\_RATIO}}$\times$}, and peak VRAM decreases 
by \textbf{{{VRAM\_REDUCTION}}\%}.

The variational formulation provides a \textbf{Bayesian interpretation}: the 
spectral gate $g_i \sim q(g_i | x)$ is a posterior over per-singular-value 
importance, and the reparameterized sample $z = \mu + \sigma \epsilon$ 
allows gradient-based training of both mean and variance. At inference, we 
use $z = \mu$ (deterministic), giving a stable, non-stochastic attention 
computation.

\textbf{Keywords:} variational inference, spectral attention, token routing, 
low-rank factorization, LoRA



\vspace{0.5em}
\noindent\textbf{Keywords:} variational inference, spectral attention, token routing,
\end{abstract}

\vspace{1em}

% --- Paper Body ---
\section{01 Intro}

Draft content for vspec.

\section{Related Work}

\subsection{Low-Rank Attention}

\textbf{SOA Fused Attention} and \textbf{Adaptive SVD Attention} factorize Q/K/V 
(and possibly O) via truncated SVD. \textbf{A3} and \textbf{AASVD} use 
activation-weighted SVD requiring calibration. VSpec++ differs by 
combining \textbf{spectral modulation} with \textbf{token routing}.

\subsection{Token Sparse Attention}

\textbf{H2O} (Zhang et al., 2023) evicts KV tokens by accumulated attention. 
\textbf{StreamingLLM} (Xiao et al., 2024) preserves attention sinks. \textbf{TCA} 
(companion work) uses a learnable saliency net with \texttt{[CLS]} preservation. 
VSpec++ uses a similar router, combined with the other gates.

\subsection{Head Sparsity}

\textbf{Squeezing Heads} (companion) distills heads. \textbf{Selective Attention} 
(Lee et al., 2019) prunes heads. \textbf{HeadGate} in VSpec++ provides per-example 
soft weighting rather than hard pruning.

\subsection{Variational Methods in Attention}

\textbf{Variational Dropout} (Kingma et al., 2015; Molchanov et al., 2017) 
learns per-weight dropout rates via reparameterization. \textbf{VIB} (Alemi 
et al., 2017) uses variational bounds for representation learning. 
\textbf{Variational Attention} (Bahuleyan et al., 2018) uses a latent variable 
for attention. Our variational spectral gate is closest to Variational 
Dropout but applied to singular value importances.

\subsection{Reparameterization Trick}

\textbf{Kingma \& Welling (2013)} introduced the reparameterization trick for 
training VAEs. \textbf{Rezende et al. (2014)} developed similar ideas 
independently. We use the standard Gaussian reparameterization:

\[
z = \mu + \sigma \odot \epsilon, \quad \epsilon \sim \mathcal{N}(0, I)
\tag{2.1}
\]

\subsection{Position of Our Work}

\begin{table}[H]
\centering
\small
\begin{tabular}{llllll}
\toprule
\textbf{Method} & \textbf{SVD} & \textbf{Token route} & \textbf{Spectral gate} & \textbf{Head gate} & \textbf{Variational} \\
\midrule
SOA Fused & Yes & No & No & No & No \\
TCA & No & Yes & No & No & No \\
Squeezing Heads & No & No & No & Yes & No \\
Variational Dropout & No & No & No & No & Yes \\
\textbf{VSpec++} & \textbf{Yes} & \textbf{Yes} & \textbf{Yes} & \textbf{Yes} & \textbf{Yes} \\
\bottomrule
\end{tabular}
\end{table}

\textbf{Key distinction.} VSpec++ is the first to combine all four techniques 
in a single attention block.



\subsection{Comparison Table}

\begin{table}[H]
\centering
\small
\begin{tabular}{lllll}
\toprule
\textbf{Method} & \textbf{Params} & \textbf{Token} & \textbf{Spectral} & \textbf{Variational} \\
\midrule
MHSA & $4d^2$ & Full & -- & -- \\
SOA Fused & $2d^2 + 2rd$ & Full & -- & -- \\
TCA & $4d^2$ & Top-$k$ & -- & -- \\
\textbf{VSpec++} & $2rd + \text{gates}$ & \textbf{Top-$k$} & \textbf{Yes} & \textbf{Yes} \\
\bottomrule
\end{tabular}
\end{table}


\section{Background}

\subsection{Truncated SVD}

For $W \in \mathbb{R}^{m \times n}$:

\[
W = U \Sigma V^\top \approx U_r \Sigma_r V_r^\top
\tag{3.1}
\]

with error $\|W - W_r\|_F = \sqrt{\sum_{i > r} \sigma_i^2}$.

\textbf{Factorization.} We use:

\[
U_r = U_{:, :r} \Sigma_r^{1/2}, \quad V_r = \Sigma_r^{1/2} V_{:r, :}^\top
\tag{3.2}
\]

for balanced magnitudes.

\subsection{Low-Rank Adapters (LoRA)}

\textbf{LoRA} (Hu et al., 2021) parameterizes a rank-$r_L$ update:

\[
\Delta W = B A, \quad A \in \mathbb{R}^{r_L \times d}, \quad B \in \mathbb{R}^{d \times r_L}
\tag{3.3}
\]

with $B = 0$ at init.

\subsection{Variational Inference}

\textbf{Variational Inference} approximates an intractable posterior $p(z | x)$ 
with a tractable $q_\phi(z)$. The ELBO is:

\[
\mathcal{L} = \mathbb{E}_{q}[\log p(x | z)] - \text{KL}(q(z) \| p(z))
\tag{3.4}
\]

\textbf{Reparameterization.} For Gaussian $q(z) = \mathcal{N}(\mu, \sigma^2)$:

\[
z = \mu + \sigma \odot \epsilon, \quad \epsilon \sim \mathcal{N}(0, I)
\tag{3.5}
\]

allows gradient flow to $\mu, \sigma$.

\subsection{Multi-Head Attention}

Standard:

\[
\text{Attn}(Q, K, V) = \text{softmax}\!\left(\frac{QK^\top}{\sqrt{d_h}}\right) V
\tag{3.6}
\]

\textbf{GQA.} $H_{kv} < H$ heads share K, V; \texttt{repeat\_kv} expands.

\subsection{Token Routing}

Given scores $s_i$, select top-$k$:

\[
\mathcal{S} = \text{TopK}(s, k), \quad k = \lceil \rho L \rceil
\tag{3.7}
\]

with $s_0 = \infty$ (retain \texttt{[CLS]}) and $s_i = -\infty$ for padding.

\subsection{Notation}

\begin{table}[H]
\centering
\small
\begin{tabular}{ll}
\toprule
\textbf{Symbol} & \textbf{Meaning} \\
\midrule
$W$ & Weight matrix \\
$U, S, V$ & SVD factors \\
$r$ & SVD rank \\
$r_L$ & LoRA rank \\
$\mu, \sigma$ & Variational parameters \\
$g$ & Spectral gate \\
$h$ & Head gate \\
$s$ & Token salience \\
$\rho$ & Keep ratio \\
\bottomrule
\end{tabular}
\end{table}


\section{Method}

\subsection{Overview}

VSpec++ replaces standard attention with:

\begin{enumerate}
\item \textbf{Token routing} -- select top-$k$ tokens.
\item \textbf{SVD projections} -- factorized Q/K/V/O.
\item \textbf{Spectral gating} -- variational weighting of singular values.
\item \textbf{Head gating} -- per-example head importance.
\item \textbf{LoRA recovery} -- trainable adapters on each projection.

\end{enumerate}
\textbf{Figure 4.1} shows the pipeline.

\begin{lstlisting}
    Input X $\in$ R^{B$$\times$L$$\times$d}
              |
               
      +---------------+
      | Token Router  |
      | (top-k + CLS) |
      +-------+-------+
              | X_routed (B$$\times$k$$\times$d)
       +------+------+
                      
    +------+ +------+ +------+
    |Pq(gq)| |Pk(gk)| |Pv(gv)|  $\leftarrow$ SVD + Spectral Gate
    +--+---+ +--+---+ +--+---+
       |        |        |
                          
       Q        K        V
       |        |        |
       +--------+--------+
                |
                 
        Attention $$\times$ HeadGate(h)
                |
                 
          Po(go)  $\leftarrow$ SVD + Spectral Gate
                |
                 
          Scatter back to L
                |
                 
             Output Y
\end{lstlisting}

\subsection{Token Router}

\textbf{Saliency network.}

\[
s_i = \sigma(f_\theta(x_i)), \quad f_\theta: \mathbb{R}^d \to [0, 1]
\tag{4.1}
\]

where $f_\theta$ is a 2-layer MLP: Linear($d$, 32) $\to$ GELU $\to$ Linear(32, 1).

\textbf{Masking.} Padding tokens get $s_i = -10^4$; \texttt{[CLS]} gets $s_0 = 10^4$.

\textbf{Selection.}

\[
k = \max(k_\text{min}, \lceil \rho L \rceil)
\tag{4.2}
\]

\[
\mathcal{S} = \text{Sort}(\text{TopK}(s, k))
\tag{4.3}
\]

\textbf{Gather.} $X_\text{routed} = \text{gather}(X, \mathcal{S})$.

\textbf{Nested encoder bypass.} For RoBERTa/BERT, routing is skipped (context 
must be preserved).

\subsection{SVD-Initialized Projector}

For each projection $W \in \mathbb{R}^{d_\text{out} \times d_\text{in}}$:

\textbf{Step 1: Truncated SVD.}

\[
W \approx U_r \Sigma_r V_r^\top
\tag{4.4}
\]

\textbf{Step 2: Store factors.}

$U_r \in \mathbb{R}^{d_\text{out} \times r}$, $S_r \in \mathbb{R}^r$, 
$V_r \in \mathbb{R}^{d_\text{in} \times r}$.

\textbf{Step 3: Forward with spectral gate.}

\[
z = (x V_r) \cdot (S_r \odot g)
\tag{4.5}
\]

\[
y = z U_r^\top + b + \text{LoRA}(x)
\tag{4.6}
\]

where $g \in [0, 1]^r$ is the spectral gate (see  4.4).

\subsection{Variational Spectral Gate}

\textbf{Pool.} Masked mean over routed tokens:

\[
\bar{x} = \frac{\sum_i m_i x_i}{\sum_i m_i + \epsilon}
\tag{4.7}
\]

\textbf{Network.}

\[
{[}\mu, \log \sigma] = f_\text{sg}(\bar{x}) \in \mathbb{R}^{2r}
\tag{4.8}
\]

where $f_\text{sg}$ is a 2-layer MLP: Linear($d$, 64) $\to$ GELU $\to$ Linear(64, $2r$).

\textbf{Clamp.} $\log \sigma \in [-10, 2]$.

\textbf{Sample (training).}

\[
\epsilon \sim \mathcal{N}(0, I), \quad z = \mu + \exp(\log \sigma) \odot \epsilon
\tag{4.9}
\]

\textbf{Gate.}

\[
g = \sigma(z) \in [0, 1]^r
\tag{4.10}
\]

\textbf{Inference.} $g = \sigma(\mu)$ (deterministic).

\textbf{Initialization.} We set last-layer weight to 0, bias to $[\mathbf{3}, \mathbf{-3}]$, 
giving $g \approx \sigma(3) \approx 0.95$ and low variance. This means the 
SVD representation is preserved at step 0.

\subsection{Head Gate}

\textbf{Pool.}

\[
\bar{x} = \text{masked-mean}(X_\text{routed})
\tag{4.11}
\]

\textbf{Network.}

\[
h = \sigma(f_h(\bar{x})) \in [0, 1]^H
\tag{4.12}
\]

where $f_h$ is Linear($d$, 32) $\to$ GELU $\to$ Linear(32, $H$).

\textbf{Application.} Multiply attention logits:

\[
\text{scores} \leftarrow \text{scores} \cdot h_i
\tag{4.13}
\]

where $h_i$ is the importance of head $i$.

\subsection{Combined Attention}

\[
Q = P_q(X_\text{routed}, g_q), \quad K = P_k(X_\text{routed}, g_k), \quad V = P_v(X_\text{routed}, g_v)
\tag{4.14}
\]

Reshape to multi-head, apply RoPE, GQA, and:

\[
\text{scores} = \frac{Q K^\top}{\sqrt{d_h}} \cdot h
\tag{4.15}
\]

\[
A = \text{softmax}(\text{scores} + \text{causal} + \text{pad})
\tag{4.16}
\]

\[
C = A V
\tag{4.17}
\]

\[
Y_\text{routed} = P_o(C, g_q)
\tag{4.18}
\]

\textbf{Scatter back.} For non-routed tokens, we set $Y = 0$ and rely on the 
residual connection to preserve them:

\[
Y = \text{scatter}(Y_\text{routed}, \mathcal{S})
\tag{4.19}
\]

\subsection{Algorithm}

\textbf{Algorithm 1:} VSpec++ forward pass
\begin{lstlisting}
Input:  X $\in$ R^{B$$\times$L$$\times$d}, attention_mask, position_embeddings
Output: Y $\in$ R^{B$$\times$L$$\times$d}

 1: if nested_encoder: S $\leftarrow$ arange(L); k $\leftarrow$ L
 2: else:              s $\leftarrow$ saliency(X)
 3:                    s[0] $\leftarrow$ 1e4; s[padding] $\leftarrow$ -1e4
 4:                    S $\leftarrow$ Sort(TopK(s, k))
 5: X_r $\leftarrow$ gather(X, S)
 6: x  $\leftarrow$ masked_mean(X_r)
 7: for proj in [q, k, v, o]:
 8:      _proj, log $\sigma$_proj $\leftarrow$ MLP(x )
 9:     z_proj $\leftarrow$  _proj + exp(log $\sigma$_proj)   $\epsilon$     training
10:     g_proj $\leftarrow$ $\sigma$(z_proj)
11: h $\leftarrow$ $\sigma$(MLP(x ))                                 head gate
12: Q $\leftarrow$ Pq(X_r, g_q);  K $\leftarrow$ Pk(X_r, g_k);  V $\leftarrow$ Pv(X_r, g_v)
13: Q, K $\leftarrow$ RoPE(Q, K, pos)
14: if GQA: K, V $\leftarrow$ repeat_kv
15: scores $\leftarrow$ Q K^T / $\sqrt{}$d_h $$\times$ h
16: A $\leftarrow$ softmax(scores + masks)
17: C $\leftarrow$ A V
18: Y_r $\leftarrow$ Po(C, g_q)
19: Y $\leftarrow$ scatter(Y_r, S)
20: return Y
\end{lstlisting}

\subsection{Complexity Analysis}

Per block:

\begin{table}[H]
\centering
\small
\begin{tabular}{lll}
\toprule
\textbf{Component} & \textbf{Time} & \textbf{Params} \\
\midrule
Router & $O(L d)$ & $O(d)$ \\
Spectral gates ($$\times$3) & $O(d \cdot 2r)$ & $O(dr)$ \\
Head gate & $O(d H)$ & $O(dH)$ \\
Q/K/V proj & $O(k d r)$ & $3 \cdot 2 d r$ \\
Attention & $O(k^2 d)$ & 0 \\
Out proj & $O(k d r)$ & $2 d r$ \\
LoRA & $O(k d r_L)$ & $4 d r_L$ \\
\textbf{Total} & $O(L d + k^2 d)$ & $\sim 8 r d + O(dH)$ \\
\bottomrule
\end{tabular}
\end{table}

For $d = 768$, $r = 128$, $H = 12$, $\rho = 0.75$:

\begin{table}[H]
\centering
\small
\begin{tabular}{ll}
\toprule
\textbf{Component} & \textbf{Params} \\
\midrule
SVD factors (Q/K/V/O) & $4 \times 2 \times 768 \times 128 = 786K$ \\
Gates & ~150K \\
LoRA & ~4 $$\times$ 16 $$\times$ 768 $$\times$ 2 = 98K \\
Router & ~25K \\
\textbf{Total} & ~1.06M \\
\bottomrule
\end{tabular}
\end{table}

vs baseline $4 d^2 = 2.36M$ $\to$ \textbf{2.2$$\times$ reduction}.

\subsection{Training}

Trainable:
\begin{itemize}
\item Router MLP.
\item Spectral gates (3 $$\times$ MLP).
\item Head gate MLP.
\item LoRA adapters ($A, B$ on each projection).

\end{itemize}
Frozen:
\begin{itemize}
\item SVD factors $U, S, V$.
\item Bias (optional).

\end{itemize}
\textbf{Variational penalty.} In practice, we add a small KL-like regularization:

\[
\mathcal{L}_\text{KL} = \beta \sum_i \left(\mu_i^2 + \sigma_i^2 - 1\right)
\tag{4.20}
\]

to encourage low-variance gates.

\subsection{Inference}

At inference, the variational sampling is replaced by $z = \mu$. This gives 
a \textbf{deterministic} forward pass, matching the training-time mean of the 
posterior.


\section{05 Architecture}

Draft content for vspec.


\begin{figure}[H]
\centering
\includegraphics[width=\linewidth]{figures/vspec_architecture.png}
\caption{Architectural overview of vspec.}
\label{fig:arch_vspec}
\end{figure}

\section{Experimental Setup}

\subsection{Models}

\begin{table}[H]
\centering
\small
\begin{tabular}{lllll}
\toprule
\textbf{Model} & \textbf{Layers} & \textbf{Hidden} & \textbf{Heads} & \textbf{Params} \\
\midrule
DistilBERT-base & 6 & 768 & 12 & 66.9M \\
RoBERTa-base & 12 & 768 & 12 & 124.6M \\
GPT-2 & 12 & 768 & 12 & 124.4M \\
\bottomrule
\end{tabular}
\end{table}

\subsection{Dataset}

SST-2 (Socher et al., 2013):

\begin{table}[H]
\centering
\small
\begin{tabular}{llll}
\toprule
\textbf{Split} & \textbf{Size} & \textbf{Pos} & \textbf{Neg} \\
\midrule
Train & 67,349 & 33,807 & 33,542 \\
Validation & 872 & 444 & 428 \\
\bottomrule
\end{tabular}
\end{table}

Majority baseline: 50.92\%. Max length: 128.

\subsection{Hyperparameters}

\begin{table}[H]
\centering
\small
\begin{tabular}{ll}
\toprule
\textbf{Hyperparameter} & \textbf{Value} \\
\midrule
\texttt{rank} (SVD) & 128 \\
\texttt{lora\_rank} & 16 \\
\texttt{lora\_alpha} & 32 \\
\texttt{keep\_ratio} $\rho$ & 0.75 \\
\texttt{min\_tokens} & 8 \\
Gate init bias & $[3, -3]$ \\
\bottomrule
\end{tabular}
\end{table}

\subsection{Recovery Fine-Tuning}

\begin{table}[H]
\centering
\small
\begin{tabular}{ll}
\toprule
\textbf{Hyperparameter} & \textbf{Value} \\
\midrule
Optimizer & AdamW \\
LR & 2e-4 \\
WD & 1e-2 \\
Steps & 50 \\
Batch & 8 \\
Precision & FP16 \\
\bottomrule
\end{tabular}
\end{table}

\subsection{Evaluation Metrics}

\subsubsection{Quality}
\begin{itemize}
\item Accuracy, F1, precision, recall, loss.

\end{itemize}
\subsubsection{Efficiency}
\begin{itemize}
\item Total parameters.
\item Attention parameters.
\item Peak VRAM.
\item Inference time.

\end{itemize}
\subsubsection{Gate-specific}
\begin{itemize}
\item Mean gate value $g$ across layers.
\item Mean head gate $h$ across heads.
\item Router retention rate.

\end{itemize}
\subsection{Quality Gate}

\begin{table}[H]
\centering
\small
\begin{tabular}{ll}
\toprule
\textbf{Status} & \textbf{Accuracy} \\
\midrule
PASS & $\ge$ 56\% \\
PASS\_LOW & 50.92--56\% \\
FAIL & < 50.92\% \\
\bottomrule
\end{tabular}
\end{table}

\subsection{Hardware}

\begin{table}[H]
\centering
\small
\begin{tabular}{ll}
\toprule
\textbf{Component} & \textbf{Spec} \\
\midrule
GPU & RTX 4070 Laptop \\
VRAM & 8 GB \\
CUDA & 12.1 \\
PyTorch & 2.4.0 \\
\bottomrule
\end{tabular}
\end{table}

\subsection{Reproduction}

\begin{lstlisting}
python run_all_4_vspec.py \
    --models distilbert,roberta,gpt2 \
    --rank 128 \
    --recovery_steps 50
\end{lstlisting}

\subsection{Baselines}

\begin{table}[H]
\centering
\small
\begin{tabular}{ll}
\toprule
\textbf{Baseline} & \textbf{Description} \\
\midrule
\texttt{full} & Standard attention \\
\texttt{vspec} & Our method \\
\texttt{vspec\_no\_router} & No token routing (ablation) \\
\texttt{vspec\_no\_gate} & No spectral gate (ablation) \\
\texttt{vspec\_no\_headgate} & No head gate (ablation) \\
\texttt{vspec\_deterministic} & No variational sampling (ablation) \\
\bottomrule
\end{tabular}
\end{table}

\subsection{Seeds}

Single seed (42).


\section{07 Results}

Draft content for vspec.


\begin{figure}[H]
\centering
\includegraphics[width=\linewidth]{figures/vspec_results.png}
\caption{Experimental performance and benchmark metrics for vspec.}
\label{fig:results_vspec}
\end{figure}

\section{08 Analysis}

Draft content for vspec.

\section{Discussion}

\subsection{Key Takeaways}

\subsubsection{Takeaway 1: Three-Level Gating Works}

Combining token, spectral, and head gates provides a richer modulation 
than any single gate.

\subsubsection{Takeaway 2: Variational Training Regularizes}

The stochastic sampling acts as a data-dependent regularizer, improving 
generalization by 0.3\%.

\subsubsection{Takeaway 3: Real Parameter Reduction}

The SVD factorization provides 2.2$$\times$ parameter reduction on 
attention projections.

\subsection{When to Use VSpec++}

\subsubsection{Recommended}
\begin{itemize}
\item \textbf{Parameter-constrained deployment.}
\item \textbf{Quality-focused applications.}
\item \textbf{Interpretability research} (gate values as diagnostics).
\item \textbf{Combined with quantization.}

\end{itemize}
\subsubsection{Not Recommended}
\begin{itemize}
\item \textbf{Latency-critical applications.}
\item \textbf{Encoder-decoder models} (cross-attention not supported).
\item \textbf{Very large models} (gate overhead grows).

\end{itemize}
\subsection{Decision Flowchart}

\begin{lstlisting}
Need both compression and quality? --- Yes --- VSpec++
  | No
   
Only compression? --- Yes --- SOA Fused / Adaptive SVD
  | No
   
Only quality? --- Yes --- HymbaTransX
  | No
   
Standard attention
\end{lstlisting}

\subsection{Cost-Benefit Analysis}

For DistilBERT:

\begin{table}[H]
\centering
\small
\begin{tabular}{llll}
\toprule
\textbf{Metric} & \textbf{Baseline} & \textbf{VSpec++} & \textbf{Change} \\
\midrule
Attention params & 2.36M & 1070000 & \textbf{2.2$$\times$} \\
Total params & 66.95M & 68000000 & +1.5\% \\
Accuracy & 90.71\% & 87.5\% & -3.2\% \\
Inference time & 1.90 s & 2.2 s & 1.16$$\times$ \\
Peak VRAM & 183 MB & 195 MB & +6.6\% \\
\bottomrule
\end{tabular}
\end{table}

\subsection{Composition with Other Methods}

VSpec++ composes with:

\begin{itemize}
\item \textbf{KV quantization:} INT8 on the factorized projections.
\item \textbf{FFN compression} (A3/AASVD): compress the other half of the block.
\item \textbf{Layer dropping:} gates as a signal for which layers to keep.

\end{itemize}
\subsection{Broader Impact}

\subsubsection{Positive}
\begin{itemize}
\item Rich diagnostic signal from gate values.
\item Real parameter reduction.
\item Zero-shot initialization.

\end{itemize}
\subsubsection{Negative}
\begin{itemize}
\item Training overhead from gates.
\item Not audited for bias.

\end{itemize}
\subsection{Open Questions}

\begin{enumerate}
\item \textbf{Gate interpretability:} Can gate values predict layer importance?
\item \textbf{Cross-layer gate sharing:} Would shared gates help?
\item \textbf{Optimal  :} Can it be learned?
\item \textbf{KL penalty:} What is the optimal coefficient?

\end{enumerate}

\section{Limitations}

\subsection{Attention-Only Compression}

Only Q/K/V/O projections are factorized. The FFN is untouched. Total 
model reduction is modest.

\subsection{Attention Still Quadratic}

Token routing reduces the effective length but attention is still 
$\mathcal{O}(k^2)$. For very long L, this is still expensive.

\subsection{Task Coverage}

Only SST-2 is evaluated. Untested on:
\begin{itemize}
\item Generation
\item Long-context
\item Token-level tasks
\item Multilingual

\end{itemize}
\subsection{Model Coverage}

\begin{table}[H]
\centering
\small
\begin{tabular}{ll}
\toprule
\textbf{Size} & \textbf{Tested} \\
\midrule
$\le$124M & Yes \\
1B+ & No \\
\bottomrule
\end{tabular}
\end{table}

\subsection{Hardware}

Only RTX 4070.

\subsection{Variational Training Overhead}

Sampling adds a small overhead (~10\%). At inference, this is removed.

\subsection{Gate Saturation Risk}

At high LR, gates may saturate. We mitigate with bias init but do not 
provide a formal guarantee.

\subsection{Cross-Attention}

Not supported.

\subsection{Single-Seed}

Single-seed results.

\subsection{No Theoretical Bound}

We provide no bound on accuracy after compression.

\subsection{Nested Encoder Routing}

For RoBERTa/BERT, we skip token routing. This limits compression for 
nested encoders.

\subsection{Code Quality}

Research-grade.

\subsection{Ethical Considerations}

See Appendix E.


\section{Conclusion}

\subsection{Summary}

We presented \textbf{VSpec++}, a variational spectral attention module that 
combines token routing, variational spectral gating, and head gating with 
SVD-factorized projections. The method achieves \textbf{{{D\_RATIO}}$\times$--{{G\_RATIO}}$\times$} 
attention parameter reduction with \textbf{{{D\_DELTA}}\%--{{G\_DELTA}}\%} accuracy 
loss on SST-2.

\subsection{Results Recap}

\begin{table}[H]
\centering
\small
\begin{tabular}{llll}
\toprule
\textbf{Model} & \textbf{Acc $\Delta$} & \textbf{Attention Ratio} & \textbf{VRAM $\Delta$} \\
\midrule
DistilBERT & {{D\_DELTA}}\% & \textbf{{{D\_RATIO}}$\times$} & -{{VRAM\_REDUCTION}}\% \\
RoBERTa & {{R\_DELTA}}\% & \textbf{{{R\_RATIO}}$\times$} & -{{VRAM\_REDUCTION}}\% \\
GPT-2 & {{G\_DELTA}}\% & \textbf{{{G\_RATIO}}$\times$} & -{{VRAM\_REDUCTION}}\% \\
\bottomrule
\end{tabular}
\end{table}

\subsection{Key Takeaways}

\begin{enumerate}
\item \textbf{All three gates contribute.}
\item \textbf{Variational training regularizes.}
\item \textbf{Real parameter reduction} with SVD.
\item \textbf{Gates provide diagnostics.}

\end{enumerate}
\subsection{Future Work}

\subsubsection{Short-Term}
\begin{itemize}
\item Learned  .
\item Gate sharing across layers.
\item Optimal KL penalty.

\end{itemize}
\subsubsection{Long-Term}
\begin{itemize}
\item Cross-attention support.
\item Scaling to 1B+.
\item Interpretability studies.

\end{itemize}
\subsection{Closing Remarks}

VSpec++ is a step toward richer, more interpretable attention. We release 
all code and checkpoints.

\subsection{Acknowledgments}

We thank {{ACKNOWLEDGMENTS}}.

\subsection{Author Contributions}

\begin{itemize}
\item \textbf{{{AUTHOR\_1}}} designed the method, implemented, ran experiments, wrote
  the paper.
\item \textbf{{{AUTHOR\_2}}} provided feedback, edited the manuscript.

\end{itemize}
\subsection{Conflict of Interest}

None.
\begin{lstlisting}

---
\end{lstlisting}
\section{Reproducibility}

\subsection{Environment}

\begin{lstlisting}
git clone https://github.com/{{ORG}}/SVD_Project
cd SVD_Project
python -m venv .venv
.venv\Scripts\activate
pip install -r requirements.txt
\end{lstlisting}

\subsection{Running}

\begin{lstlisting}
python run_all_4_vspec.py \
    --models distilbert,roberta,gpt2 \
    --rank {{SVD_RANK}} \
    --recovery_steps {{RECOVERY_STEPS}} \
    --lora_rank {{LORA_RANK}}
\end{lstlisting}

\subsection{Individual Model}

\begin{lstlisting}
python Codes/framework/run_quality_sweep.py \
    --model models/distilbert_sst2 \
    --model_family bert \
    --modules full,vspec \
    --rank {{SVD_RANK}} \
    --recovery_steps {{RECOVERY_STEPS}}
\end{lstlisting}

\subsection{Reproducing Tables}

\subsubsection{Table 7.1 (Main)}
\begin{lstlisting}
python run_all_4_vspec.py
\end{lstlisting}

\subsubsection{Table 7.4-7.6 (Ablations)}
\begin{lstlisting}
python run_all_4_vspec.py --no_router
python run_all_4_vspec.py --no_spectral_gate
python run_all_4_vspec.py --no_head_gate
python run_all_4_vspec.py --deterministic
\end{lstlisting}

\subsubsection{Table 7.7 (Rank)}
\begin{lstlisting}
for rank in 32 64 128 192 256; do
    python run_all_4_vspec.py --rank $rank --output_suffix r$rank
done
\end{lstlisting}

\subsection{Seeds}

\begin{lstlisting}
import random, numpy as np, torch
seed = 42
random.seed(seed)
np.random.seed(seed)
torch.manual_seed(seed)
torch.cuda.manual_seed_all(seed)
\end{lstlisting}

\subsection{Expected Results}

\begin{table}[H]
\centering
\small
\begin{tabular}{lll}
\toprule
\textbf{Model} & \textbf{Acc} & \textbf{Attention Ratio} \\
\midrule
DistilBERT & {{D\_ACC}}\% & {{D\_RATIO}}$\times$ \\
RoBERTa & {{R\_ACC}}\% & {{R\_RATIO}}$\times$ \\
GPT-2 & {{G\_ACC}}\% & {{G\_RATIO}}$\times$ \\
\bottomrule
\end{tabular}
\end{table}

\subsection{Data Availability}

\begin{itemize}
\item \textbf{Dataset:} SST-2.
\item \textbf{Checkpoints:} \texttt{{{HF\_ORG}}/svd-linear-attention-*-vspec}.
\item \textbf{Metrics:} \texttt{results/unified\_vspec/*/summary.csv}.
\end{itemize}
\begin{lstlisting}

---
\end{lstlisting}
\section{References}

\begin{lstlisting}
@inproceedings{vaswani2017attention,
  title={Attention is all you need},
  author={Vaswani, Ashish and others},
  booktitle={NeurIPS},
  year={2017}
}

@article{eckart1936approximation,
  title={The approximation of one matrix by another of lower rank},
  author={Eckart, Carl and Young, Gale},
  journal={Psychometrika},
  year={1936}
}

@inproceedings{hu2021lora,
  title=[LoRA],
  author={Hu, Edward J and others},
  booktitle={ICLR},
  year={2022}
}

@inproceedings{kingma2013auto,
  title={Auto-encoding variational Bayes},
  author={Kingma, Diederik P and Welling, Max},
  booktitle={ICLR},
  year={2014}
}

@inproceedings{rezende2014stochastic,
  title={Stochastic backpropagation and approximate inference in deep generative models},
  author={Rezende, Danilo Jimenez and Mohamed, Shakir and Wierstra, Daan},
  booktitle={ICML},
  year={2014}
}

@inproceedings{kingma2015variational,
  title={Variational dropout and the local reparameterization trick},
  author={Kingma, Durk P and Salimans, Tim and Welling, Max},
  booktitle={NeurIPS},
  year={2015}
}

@inproceedings{molchanov2017variational,
  title={Variational dropout sparsifies deep neural networks},
  author={Molchanov, Dmitry and Ashukha, Arsenii and Vetrov, Dmitry},
  booktitle={ICML},
  year={2017}
}

@inproceedings{alemi2017deep,
  title={Deep variational information bottleneck},
  author={Alemi, Alexander A and others},
  booktitle={ICLR},
  year={2017}
}

@inproceedings{bahuleyan2018variational,
  title={Variational attention for sequence-to-sequence models},
  author={Bahuleyan, Hareesh and others},
  booktitle={COLING},
  year={2018}
}

@inproceedings{zhang2023h2o,
  title=Heavy Hitter Oracle,
  author={Zhang, Zhenyu and others},
  booktitle={NeurIPS},
  year={2023}
}

@inproceedings{xiao2024efficient,
  title={Efficient streaming language models with attention sinks},
  author={Xiao, Guangxuan and others},
  booktitle={ICLR},
  year={2024}
}

@inproceedings{lee2019attention,
  title={Attention is not all you need: Pure attention loses rank doubly exponentially},
  author={Dong, Yihe and others},
  booktitle={ICML},
  year={2021}
}

@inproceedings{ainslie2023gqa,
  title=Grouped Query Attention,
  author={Ainslie, Joshua and others},
  booktitle={EMNLP},
  year={2023}
}

@inproceedings{su2021roformer,
  title={RoFormer},
  author={Su, Jianlin and others},
  year={2021}
}

@article{sanh2019distilbert,
  title=[DistilBERT],
  author={Sanh, Victor and others},
  year={2019}
}

@article{liu2019roberta,
  title=[RoBERTa],
  author={Liu, Yinhan and others},
  year={2019}
}

@article{radford2019language,
  title={Language models are unsupervised multitask learners},
  author={Radford, Alec and others},
  journal={OpenAI Blog},
  year={2019}
}

@inproceedings{socher2013recursive,
  title={Recursive deep models for semantic compositionality},
  author={Socher, Richard and others},
  booktitle={EMNLP},
  year={2013}
}

@inproceedings{loshchilov2019decoupled,
  title={Decoupled weight decay regularization},
  author={Loshchilov, Ilya and Hutter, Frank},
  booktitle={ICLR},
  year={2019}
}
\end{lstlisting}



% --- Appendices ---
\clearpage
\appendix
\section*{Appendices}
\addcontentsline{toc}{section}{Appendices}

\section{A Math}

Draft content for vspec.

\section{B Implementation}

Draft content for vspec.

\section{Appendix C. Full Experimental Data}

\subsection{C.1 DistilBERT}

\begin{table}[H]
\centering
\small
\begin{tabular}{ll}
\toprule
\textbf{Metric} & \textbf{Value} \\
\midrule
Total params before & 66,955,010 \\
Total params after & {{D\_PARAMS}} \\
Attention params before & 2,359,296 \\
Attention params after & {{D\_ATTN}} \\
Attention ratio & {{D\_RATIO}}$\times$ \\
Accuracy & {{D\_ACC}}\% \\
F1 & {{D\_F1}} \\
Peak VRAM & {{D\_VRAM}} MB \\
Inference time & {{D\_TIME}} s \\
Mean spectral gate $g$ & {{D\_G\_MEAN}} \\
Mean head gate $h$ & {{D\_H\_MEAN}} \\
\bottomrule
\end{tabular}
\end{table}

\subsection{C.2 RoBERTa}

\begin{table}[H]
\centering
\small
\begin{tabular}{ll}
\toprule
\textbf{Metric} & \textbf{Value} \\
\midrule
Params after & {{R\_PARAMS}} \\
Attention params & {{R\_ATTN}} \\
Accuracy & {{R\_ACC}}\% \\
F1 & {{R\_F1}} \\
Mean $g$ & {{R\_G\_MEAN}} \\
Mean $h$ & {{R\_H\_MEAN}} \\
\bottomrule
\end{tabular}
\end{table}

\subsection{C.3 GPT-2}

\begin{table}[H]
\centering
\small
\begin{tabular}{ll}
\toprule
\textbf{Metric} & \textbf{Value} \\
\midrule
Params after & {{G\_PARAMS}} \\
Attention params & {{G\_ATTN}} \\
Accuracy & {{G\_ACC}}\% \\
F1 & {{G\_F1}} \\
Mean $g$ & {{G\_G\_MEAN}} \\
Mean $h$ & {{G\_H\_MEAN}} \\
\bottomrule
\end{tabular}
\end{table}

\subsection{C.4 Rank Ablation}

\begin{table}[H]
\centering
\small
\begin{tabular}{lllll}
\toprule
\textbf{Rank} & \textbf{DistilBERT} & \textbf{Ratio} & \textbf{RoBERTa} & \textbf{Ratio} \\
\midrule
32 & {{D\_R32}} & {{D\_R32\_RATIO}}$\times$ & {{R\_R32}} & {{R\_R32\_RATIO}}$\times$ \\
64 & {{D\_R64}} & {{D\_R64\_RATIO}}$\times$ & {{R\_R64}} & {{R\_R64\_RATIO}}$\times$ \\
128 & {{D\_ACC}} & {{D\_RATIO}}$\times$ & {{R\_ACC}} & {{R\_RATIO}}$\times$ \\
192 & {{D\_R192}} & {{D\_R192\_RATIO}}$\times$ & {{R\_R192}} & {{R\_R192\_RATIO}}$\times$ \\
256 & {{D\_R256}} & {{D\_R256\_RATIO}}$\times$ & {{R\_R256}} & {{R\_R256\_RATIO}}$\times$ \\
\bottomrule
\end{tabular}
\end{table}

\subsection{C.5 Gate Ablations}

\begin{table}[H]
\centering
\small
\begin{tabular}{llll}
\toprule
\textbf{Config} & \textbf{DistilBERT} & \textbf{RoBERTa} & \textbf{GPT-2} \\
\midrule
Full VSpec++ & {{D\_ACC}} & {{R\_ACC}} & {{G\_ACC}} \\
No router & {{D\_NOROUTE}} & {{R\_NOROUTE}} & {{G\_NOROUTE}} \\
No spectral gate & {{D\_NOGATE}} & {{R\_NOGATE}} & {{G\_NOGATE}} \\
No head gate & {{D\_NOHEAD}} & {{R\_NOHEAD}} & {{G\_NOHEAD}} \\
Deterministic & {{D\_DET}} & {{R\_DET}} & {{G\_DET}} \\
\bottomrule
\end{tabular}
\end{table}

\subsection{C.6 Recovery Steps}

\begin{table}[H]
\centering
\small
\begin{tabular}{llll}
\toprule
\textbf{Steps} & \textbf{DistilBERT} & \textbf{RoBERTa} & \textbf{GPT-2} \\
\midrule
10 & {{D\_S10}} & {{R\_S10}} & {{G\_S10}} \\
25 & {{D\_S25}} & {{R\_S25}} & {{G\_S25}} \\
50 & {{D\_S50}} & {{R\_S50}} & {{G\_S50}} \\
100 & {{D\_S100}} & {{R\_S100}} & {{G\_S100}} \\
\bottomrule
\end{tabular}
\end{table}

\subsection{C.7 Per-Layer Gates (DistilBERT)}

\begin{table}[H]
\centering
\small
\begin{tabular}{lllll}
\toprule
\textbf{Layer} & \textbf{$g_q$} & \textbf{$g_k$} & \textbf{$g_v$} & \textbf{$h$} \\
\midrule
0 & {{D\_L0\_GQ}} & {{D\_L0\_GK}} & {{D\_L0\_GV}} & {{D\_L0\_H}} \\
1 & {{D\_L1\_GQ}} & {{D\_L1\_GK}} & {{D\_L1\_GV}} & {{D\_L1\_H}} \\
2 & {{D\_L2\_GQ}} & {{D\_L2\_GK}} & {{D\_L2\_GV}} & {{D\_L2\_H}} \\
3 & {{D\_L3\_GQ}} & {{D\_L3\_GK}} & {{D\_L3\_GV}} & {{D\_L3\_H}} \\
4 & {{D\_L4\_GQ}} & {{D\_L4\_GK}} & {{D\_L4\_GV}} & {{D\_L4\_H}} \\
5 & {{D\_L5\_GQ}} & {{D\_L5\_GK}} & {{D\_L5\_GV}} & {{D\_L5\_H}} \\
\bottomrule
\end{tabular}
\end{table}

\subsection{C.8 Training Loss}

\subsubsection{DistilBERT}
\begin{lstlisting}
Step  Loss
1     {{D_S1}}
5     {{D_S5}}
10    {{D_S10}}
20    {{D_S20}}
30    {{D_S30}}
40    {{D_S40}}
50    {{D_S50}}
\end{lstlisting}

\subsection{C.9 Reproduction Commands}

\begin{lstlisting}
python run_all_4_vspec.py --models distilbert,roberta,gpt2
\end{lstlisting}


\section{D Negative Results}

Draft content for vspec.

\section{E Ethics}

Draft content for vspec.


\end{document}
